\section{软硬耦合分析}
微分算子 $\mathcal{L}$ 在由张成函数空间 $\mathbb{V}^h$ 的基函数 $\varphi^h_i$ 离散化后，得到的线性系统为 $\mathbf{K}^h = \int_{\Omega} \varphi \mathcal{L} \varphi^T dX$。

刚度矩阵 $\mathbf{K}$ 的病态通常源于两种情形。首先是坏形状单元引起的局部问题，即基函数 $\varphi$ 的位置不当导致其支撑区域的交集产生病态。这类单元常见于复杂的几何边界和特定的网格约束中。其次是全局问题，即算子 $\mathcal{L}$ 本身具有高度非均匀的系数和多尺度行为。因此，只要离散化足够细密，即使不存在坏形状单元，$\mathbf{K}$ 也会具有很宽的光谱。这种不良的条件数被视为一种全局问题。

针对后一种情形，已有研究提出了算子自适应可精细化基函数~\cite{chen2019material,Budninskiy_2019_wavelet}，用于构造一组相对于算子正交且在空间和特征空间中均具有局部性的基函数。以两层层级结构为例，小波Galerkin方法将原始函数空间 $\mathbb{V}^h$ 分解为两个子空间：
\begin{equation}
  \mathbb{V}^{h} = \mathbb{V}^H \oplus_{\mathcal{L}} \mathbb{W}^H,
\end{equation}
其中，$\varphi^H_i = \sum_{l} \mathbb{C}^H_{il} \varphi^h_l$ 张成函数空间 $\mathbb{V}^H$，$\psi^H_i = \sum_l \mathbf{W}^H_{il} \varphi^h_l$ 张成函数空间 $\mathbb{W}^H$。文献~\cite{chen_multiscale_2021}指出，在进行适当排序后，可以通过 Choleksy 分解计算这些层级基函数和小波函数。因此，原始细层级的刚度矩阵可以通过基变换 $\mathbf{P}$ 块对角化为 $\mathbb{K}^h = \operatorname{diag}(\mathbb{B}, \mathbb{K}^H)$。由于特征空间局部性的存在，$\mathbb{K}^H$ 或 $\mathbb{B}$ 的条件数都比 $\mathbf{K}^h$ 更小。可以直接分别求解它们，也可以使用不完全 Cholesky 分解作为预条件子，这两种方法都能得到条件数更优的线性系统。

对于前一种情形，上述方法同样适用，这在文献~\cite{chen_multiscale_2021}的开题案例中有所体现。然而，这种全局方法的计算开销非常大。此外，经过多年网格质量优化研究的发展，坏形状单元通常不会遍布整个区域。换言之，这些劣质单元往往是局部的且不可避免的（因为可能已经应用了各种重网格化技术）。因此，本研究的目标是开发一种轻量且高效的方法，专门解决这一局部的特殊问题。
